\nonstopmode{}
\documentclass[a4paper]{book}
\usepackage[times,inconsolata,hyper]{Rd}
\usepackage{makeidx}
\makeatletter\@ifl@t@r\fmtversion{2018/04/01}{}{\usepackage[utf8]{inputenc}}\makeatother
% \usepackage{graphicx} % @USE GRAPHICX@
\makeindex{}
\begin{document}
\chapter*{}
\begin{center}
{\textbf{\huge Package `directsrr'}}
\par\bigskip{\large \today}
\end{center}
\ifthenelse{\boolean{Rd@use@hyper}}{\hypersetup{pdftitle = {directsrr: Directly Standardised Rates, Rate Ratios, and Rate Differences}}}{}
\ifthenelse{\boolean{Rd@use@hyper}}{\hypersetup{pdfauthor = {Nazrul Islam}}}{}
\begin{description}
\raggedright{}
\item[Title]\AsIs{Directly Standardised Rates, Rate Ratios, and Rate Differences}
\item[Version]\AsIs{0.1.0}
\item[Description]\AsIs{Computes directly standardised rates (DSR) with confidence
intervals using the gamma-based method of Tiwari, Clegg, and Zou (2006)
<}\Rhref{https://doi.org/10.1191/1740774506cn143oa}{doi:10.1191/1740774506cn143oa}\AsIs{>, a refinement of Fay and Feuer (1997).
Optional Dobson et al. (1991) intervals are also provided
(<}\Rhref{https://doi.org/10.1002/sim.4780100317}{doi:10.1002/sim.4780100317}\AsIs{>). Between-group
comparisons include standardised rate ratios (SRR), with F-distribution
confidence intervals, and standardised rate differences (SRD), with
normal-approximation confidence intervals based on the Tiwari variance.
The package is adapted from the Stata command distrate (Consonni et al.,
2012, Stata Journal 12(4):688--701), but uses a standard-population
column in the main data frame and flexible column naming.}
\item[License]\AsIs{GPL-3}
\item[Encoding]\AsIs{UTF-8}
\item[Roxygen]\AsIs{list(markdown = TRUE)}
\item[RoxygenNote]\AsIs{7.3.3}
\item[Depends]\AsIs{R (>= 4.0.0)}
\item[Suggests]\AsIs{testthat (>= 3.0.0), knitr, rmarkdown}
\item[Config/testthat/edition]\AsIs{3}
\item[LazyData]\AsIs{true}
\item[VignetteBuilder]\AsIs{knitr}
\item[URL]\AsIs{}\url{https://github.com/Naz-EpiStats/DirectSRR}\AsIs{}
\item[BugReports]\AsIs{}\url{https://github.com/Naz-EpiStats/DirectSRR/issues}\AsIs{}
\item[NeedsCompilation]\AsIs{no}
\item[Author]\AsIs{Nazrul Islam [aut, cre] (ORCID:
<}\url{https://orcid.org/0000-0003-3982-4325}\AsIs{>)}
\item[Maintainer]\AsIs{Nazrul Islam }\email{naz@mail.harvard.edu}\AsIs{}
\end{description}
\Rdcontents{Contents}
\HeaderA{directsrr}{Directly Standardised Rates, Rate Ratios, and Rate Differences}{directsrr}
%
\begin{Description}
Calculates directly age-standardised rates (DSR) and supports pairwise
comparison of groups via standardised rate ratios (SRR) and standardised
rate differences (SRD).

Confidence intervals for the DSR use the gamma-distribution method of
Tiwari, Clegg, and Zou (2006), a refinement of Fay and Feuer (1997) that
gives empirical coverage closer to the nominal level for rare outcomes.
The Dobson et al. (1991) interval can also be requested. SRR confidence
limits use the F-distribution approach from Tiwari et al.; SRD limits use
a normal approximation based on the same Tiwari variance.

An R adaptation of the Stata command \code{distrate} (Consonni et al.
2012), with the standard population as a column in \code{data} rather than
a separate file, flexible column naming, and no run-time dependencies
beyond base R.
\end{Description}
%
\begin{Usage}
\begin{verbatim}
directsrr(
  data,
  event = "events",
  pop = "population",
  stdpop = "stdpop",
  standstrata = "age_group",
  by = NULL,
  ref = NULL,
  ref_index = NULL,
  mult = 1,
  level = 95,
  dobson = TRUE,
  check_pop = FALSE,
  digits_rate = 2,
  digits_ratio = 4,
  digits_count = 0,
  notify = c("warn", "message", "silent"),
  verbose = TRUE
)
\end{verbatim}
\end{Usage}
%
\begin{Arguments}
\begin{ldescription}
\item[\code{data}] A data frame in aggregate form. Multiple rows per stratum-group
combination are summed before any calculation.

\item[\code{event}] Column name for event counts. Must be non-negative.
Non-integer values are accepted; when \code{dobson = TRUE} the function
warns that the Dobson CI may be unreliable. Default \code{"events"}.

\item[\code{pop}] Column name for person-time or population size (denominator).
Must be strictly positive. Default \code{"population"}.

\item[\code{stdpop}] Column name for the standard-population weight. Must be
constant within each stratum and non-negative in every row; the total
weight sum within each group must be strictly positive. Absolute counts
and proportions are both accepted; values are normalised within each group.
Default \code{"stdpop"}.

\item[\code{standstrata}] Character vector of column name(s) defining the
standardisation strata. Default \code{"age\_group"}.

\item[\code{by}] Character vector of grouping column name(s), or \code{NULL} for
a single pooled rate. Default \code{NULL}.

\item[\code{ref}] Reference group for SRR and SRD, specified as a \strong{value} to
match against the \code{by} column(s). Accepts:
\begin{itemize}

\item{} A \strong{character or numeric scalar} matched against the single
\code{by} column, e.g. \code{ref = "Control"} or \code{ref = 2}.
\item{} A \strong{named list} for multi-column \code{by}, supplying a
value for every column, e.g.
\code{ref = list(arm = "Placebo", centre = "London")}.

\end{itemize}

Supply either \code{ref} or \code{ref\_index}, not both. When both are
\code{NULL} and \code{by} is not \code{NULL}, the first row of the
output table is used as reference. Ignored when \code{by = NULL}.
Default \code{NULL}.

\item[\code{ref\_index}] Reference group for SRR and SRD, specified as a
\strong{row position} in the sorted output table. A single positive whole
number (e.g. \code{1}, \code{2}, \code{3} — no \code{L} suffix
required). Supply either \code{ref} or \code{ref\_index}, not both.
Ignored when \code{by = NULL}. Default \code{NULL}.

\item[\code{mult}] Rate multiplier. Default \code{1}.

\item[\code{level}] Confidence level as a percentage. Default \code{95}.

\item[\code{dobson}] Logical. Include Dobson et al. (1991) limits? Default \code{TRUE}.

\item[\code{check\_pop}] Logical. Error if \code{event > pop} in any row? Appropriate
for binomial-style data; leave \code{FALSE} for person-time denominators.
Default \code{FALSE}.

\item[\code{digits\_rate}] Decimal places for rate and SRD columns. Default \code{2}.

\item[\code{digits\_ratio}] Decimal places for SRR columns. Default \code{4}.

\item[\code{digits\_count}] Decimal places for count columns. Default \code{0}.

\item[\code{notify}] Character. Controls how the function signals that Dobson
limits are unavailable for zero-event groups (when \code{dobson = TRUE}).
\begin{itemize}

\item{} \code{"warn"} (default) — issues a \code{warning()}, catchable
with \code{tryCatch()}. Recommended because \code{ub\_dob = NA}
changes how the interval should be interpreted.
\item{} \code{"message"} — emits a \code{message()}, suppressible with
\code{suppressMessages()}.
\item{} \code{"silent"} — no output. Useful in loops or simulations.

\end{itemize}

Has no effect when there are no zero-event groups or \code{dobson = FALSE}.

\item[\code{verbose}] Print a summary table? Default \code{TRUE}.
\end{ldescription}
\end{Arguments}
%
\begin{Value}
An invisible data frame with one row per group containing:
\code{events}, \code{N}, \code{crude}, \code{rateadj},
\code{lb\_gam}/\code{ub\_gam}, \code{se\_gam},
\code{lb\_dob}/\code{ub\_dob} (if \code{dobson = TRUE}),
\code{srr}/\code{lb\_srr}/\code{ub\_srr},
\code{srd}/\code{lb\_srd}/\code{ub\_srd}
(last six only when \code{by} is not \code{NULL}).
Rate and SRD columns are on the multiplied scale; SRR is dimensionless.
\end{Value}
%
\begin{References}
Consonni D, Coviello E, Buzzoni C, Mensi C (2012). A command to calculate
age-standardized rates with efficient interval estimation.
\emph{Stata Journal}, \strong{12}(4), 688--701.
\Rhref{https://doi.org/10.1177/1536867X1201200408}{doi:10.1177\slash{}1536867X1201200408}

Tiwari RC, Clegg LX, Zou Z (2006). Efficient interval estimation for
age-adjusted cancer rates. \emph{Statistical Methods in Medical Research},
\strong{15}(6), 547--569. \Rhref{https://doi.org/10.1191/1740774506cn143oa}{doi:10.1191\slash{}1740774506cn143oa}

Fay MP, Feuer EJ (1997). Confidence intervals for directly standardized
rates: A method based on the gamma distribution.
\emph{Statistics in Medicine}, \strong{16}(7), 791--801.

Dobson AJ, Kuulasmaa K, Eberle E, Scherer J (1991). Confidence intervals
for weighted sums of Poisson parameters.
\emph{Statistics in Medicine}, \strong{10}(3), 457--462.
\Rhref{https://doi.org/10.1002/sim.4780100317}{doi:10.1002\slash{}sim.4780100317}
\end{References}
%
\begin{Examples}
\begin{ExampleCode}
data(directsrr_ExampleData)

# Reference group identified by label
directsrr(
  data        = directsrr_ExampleData,
  event       = "events",
  pop         = "population",
  stdpop      = "stdpop",
  standstrata = "age_group",
  by          = "group",
  ref         = "A",
  mult        = 100000,
  digits_rate  = 2,
  digits_ratio = 4
)

# Reference group identified by row position (no L suffix needed)
directsrr(
  data        = directsrr_ExampleData,
  event       = "events",
  pop         = "population",
  stdpop      = "stdpop",
  standstrata = "age_group",
  by          = "group",
  ref_index   = 1,
  mult        = 100000
)

# Single pooled rate, no comparison
directsrr(
  data        = directsrr_ExampleData,
  event       = "events",
  pop         = "population",
  stdpop      = "stdpop",
  standstrata = "age_group",
  mult        = 100000
)

\end{ExampleCode}
\end{Examples}
\HeaderA{directsrr\_ExampleData}{Example dataset}{directsrr.Rul.ExampleData}
\keyword{datasets}{directsrr\_ExampleData}
%
\begin{Description}
Simulated aggregate event data for three study groups across three age
strata, intended for use in package examples and tests.
\end{Description}
%
\begin{Usage}
\begin{verbatim}
directsrr_ExampleData
\end{verbatim}
\end{Usage}
%
\begin{Format}
A data frame with 9 rows and 5 columns:
\begin{description}

\item[\code{age\_group}] Age stratum: \code{"18-39"}, \code{"40-64"}, or \code{"65+"}.
\item[\code{group}] Study group: \code{"A"}, \code{"B"}, or \code{"C"}.
\item[\code{events}] Integer event count for this stratum and group.
\item[\code{population}] Person-time or population size.
\item[\code{stdpop}] Standard population weight for the age stratum
(proportions: 0.4, 0.4, 0.2).

\end{description}

\end{Format}
%
\begin{Source}
Simulated; no real patients.
\end{Source}
%
\begin{Examples}
\begin{ExampleCode}
data(directsrr_ExampleData)
head(directsrr_ExampleData)
\end{ExampleCode}
\end{Examples}
\printindex{}
\end{document}
